\documentclass[10pt,a4paper]{amsart}
\usepackage[margin=19mm]{geometry}
\usepackage{amsmath,amssymb,mathtools}
\usepackage[scr=rsfs,cal=euler]{mathalfa}
\usepackage{xfrac}
\usepackage[unicode,hidelinks,bookmarks=false]{hyperref}
\newcommand{\ie}{i.e.\ }
\newcommand{\cf}{cf.\ }
\newcommand{\resp}{respectively\ }
\newcommand{\Subsection}{\subsection}
\providecommand{\nobreakdash}{\nobreak\hbox{-}}
\setlength{\emergencystretch}{2em}
\multlinegap0pt
\usepackage{amssymb}
\usepackage{float}
\newtheorem{lem}{Lemma}
\newtheorem{thm}{Theorem}
\title{\Huge On Clique Graphs and Clique Regular Graphs \\[1em] \Large Honors Capstone Project}
\author{Connor Phillips}
\date{Source: arXiv:2605.22867v1 (CC BY 4.0)}
\begin{document}
\maketitle

\noindent\emph{Source and licence.} \Huge On Clique Graphs and Clique Regular Graphs \\[1em] \Large Honors Capstone Project. Version arXiv:2605.22867v1 (CC BY 4.0), distributed by arXiv under the Creative Commons Attribution 4.0 International licence, http://creativecommons.org/licenses/by/4.0/, as recorded in the arXiv licence metadata for this exact version and shown on the abstract page https://arxiv.org/abs/2605.22867v1. Full licence notice: \texttt{licenses/arxiv-2605.22867-CC-BY-4.0.txt}.\par\smallskip

\noindent\textit{Editorial scope.} Counted unit: the article's thm thm:eigen (source0.tex lines 499--503). The article's complete original proof of it is delivered with it (source0.tex lines 504--524, one \texttt{proof} environment of 21 lines). No mathematics is written, repaired or re-derived: the statement and the proof are the article's own lines, in the article's own notation, and no retained line is substituted or re-flowed.  Retained uncounted as the source-local context the proof needs: lines 488--492.  Nothing was omitted from any retained span, so the package carries no omission record.  No retained line contains a citation, so the wrapper carries no bibliography and no bibliography entry.  No retained line contains a figure environment, an image-inclusion command or drawing code, so no image is delivered and the package declares no asset dependency.  Editorial formatting only, in addition: an AMS \texttt{amsart} preamble that restates the article's own declarations and macros used by the retained lines, namely the environments \texttt{lem}, \texttt{thm} on separate counters, together with the standard packages those lines need.  The retained environments print in this document as Lemma 1, Theorem 1 in order of appearance, whereas the article prints the same items with the numbering of its own full document.  The complete native source of the article as distributed in the arXiv e-print for this exact version is bundled unchanged as \texttt{sources/201165/source0.tex}.
\begin{lem} \label{lem:blockmat}
Suppose that $\Gamma$ is $\omega$-clique regular with $\omega$-clique incidence matrix $R$ and degree matrix $D$. Then:\vspace{0.2cm} \\
           \hspace*{0.5cm} (1) $R^\top R=A_C + \omega I_m$ and  \\
           \hspace*{0.5cm} (2) $RR^\top  = A + \frac{1}{\omega -1}D$. \vspace{0.2cm}
\end{lem}

\begin{thm} \label{thm:eigen}
    If $\Gamma$ is $k$-regular and $\omega$-clique regular, then 
    \[p(C_\omega(\Gamma); \lambda)=(\lambda+\omega)^{m-n}p\left(\Gamma;\lambda+\omega-\frac{k}{\omega - 1}\right)\]
    where $m=\frac{nk}{\omega(\omega-1)}$.
\end{thm}

\begin{proof}
    Define two square block matrices with $n+m$ rows and columns as follows with $R$ the $\omega$-clique incidence matrix of $\Gamma$,
    \[U = \begin{bmatrix}(\lambda+\omega) I_n & -R \\0 & I_m \\\end{bmatrix}, \hspace{1cm}
    V = \begin{bmatrix}I_n & R \\R^\top  & (\lambda+\omega) I_m \\\end{bmatrix}.\]
Then, 
\[UV = \begin{bmatrix}(\lambda+\omega) I_n - RR^\top  & 0 \\R^\top  & (\lambda+\omega) I_m \\\end{bmatrix}, 
    VU = \begin{bmatrix}(\lambda+\omega) I_n & 0 \\(\lambda+\omega) R^\top  & (\lambda+\omega) I_m - R^\top R \\\end{bmatrix}.\]
Since $\text{det}(UV)=\text{det}(VU)$, it follows that
\begin{align*}
    (\lambda+\omega)^m \text{det}((\lambda+\omega) I_n - RR^\top ) & =(\lambda+\omega)^n \text{det}((\lambda+\omega) I_m - R^\top R) \\ 
    (\lambda+\omega)^{m-n} \text{det}((\lambda+\omega) I_n - RR^\top ) & =\text{det}((\lambda+\omega) I_m - R^\top R).
\end{align*}
Using the above equality and Lemma \ref{lem:blockmat} we get,
\begin{align*}
    p(C_\omega(\Gamma);\lambda) & = \text{det}(\lambda I_m - A_C) \\ 
    & = \text{det}((\lambda + \omega)I_m - R^\top R)  \\
    & = (\lambda + \omega)^{m-n}\text{det}((\lambda + \omega)I_n - RR^\top ) \\
    & = (\lambda + \omega)^{m-n}\text{det}\left(\left(\lambda + \omega - \frac{k}{\omega -1}\right)I_n - A\right)\\
    & = (\lambda + \omega)^{m-n}p\left(\Gamma ; \lambda + \omega - \frac{k}{\omega -1}\right).
\end{align*}
\end{proof}

\end{document}
